% Glossary of nuclear terms (from the radioactivity source, glossary.tex)
\volumeopening{Glossary of Nuclear Terms}

Every technical term used in this volume, in alphabetical order, with the chapter where it is developed.
The notation is that of the chapters: $A$ is the mass number, $Z$ the atomic number and $A-Z$ the
neutron number; $N$ or $N(t)$ is always the number of undecayed nuclei.

{\small\emergencystretch=1.5em
\begin{description}[leftmargin=0pt,itemsep=2pt,font=\bfseries]
\item[absorbed dose ($D$)] Energy deposited by ionising radiation per unit mass of material, $D=E/m$;
  unit gray (\si{\gray}). (\chref{ch:nucapps})
\item[activation] Making a radioactive isotope by bombarding a target, often with neutrons, at a constant
  production rate $\Prod$. (\chref{ch:decaylaw})
\item[activity ($\Act$)] Number of decays per second in a sample, $\Act=\lambda N$; unit becquerel
  (\si{\becquerel}). (\chref{ch:decaylaw})
\item[alpha particle ($\alpha$)] A helium-4 nucleus (two protons and two neutrons, charge $+2e$), emitted in
  alpha decay. (\chref{ch:decay})
\item[annihilation] A particle meeting its antiparticle: both vanish and their energy appears as photons,
  for example $e^-+e^+\to2\gamma$. (\chref{ch:decay})
\item[antineutrino ($\bar\nu_e$)] The antiparticle of the electron neutrino: neutral, nearly massless,
  lepton number $-1$; emitted together with an electron in $\beta^-$ decay. (\chref{ch:decay})
\item[antiparticle, antimatter] For each particle, a partner of identical mass but opposite charge and
  opposite lepton number. Antiparticles collectively are antimatter. (\chref{ch:decay})
\item[atomic mass unit (\si{\amu})] One-twelfth of the mass of a neutral carbon-12 atom,
  \SI{1.66054e-27}{kg}; $1\,\si{\amu}\times c^2=\SI{931.494}{\MeV}$. (\chref{ch:nucleus})
\item[atomic number ($Z$)] The number of protons in a nucleus; it fixes the chemical element.
  (\chref{ch:nucleus})
\item[background] Counts recorded by a detector with no source present, from cosmic rays and natural
  radioactivity; it must be subtracted before a decay curve is analysed. (\chref{ch:decaylaw})
\item[becquerel (\si{\becquerel}), curie (\si{\curie})] Units of activity: $\SI{1}{\becquerel}$ is one decay
  per second; $\SI{1}{\curie}=\SI{3.7e10}{\becquerel}$, roughly the activity of one gram of radium-226.
  (\chref{ch:decaylaw})
\item[beta decay ($\beta^-$, $\beta^+$)] A nucleus turns a neutron into a proton, emitting an electron and an
  antineutrino ($\beta^-$), or a proton into a neutron, emitting a positron and a neutrino ($\beta^+$);
  $A$ is unchanged. (\chref{ch:decay})
\item[binding energy ($\Eb$)] The energy needed to separate a nucleus completely into free nucleons; it equals
  the mass defect times $c^2$. Its value per nucleon, $\Eb/A$, peaks near iron. (\chref{ch:nucleus})
\item[biological half-life ($T_\text{bio}$)] The time for the body to eliminate half of a substance by
  excretion alone. (\chref{ch:nucapps})
\item[branching ratio ($f_i$)] The fraction of decays that proceed by a particular channel $i$ when several are
  possible; $f_i=\lambda_i/\lambda$. (\chref{ch:decaylaw})
\item[chain reaction] Fissions releasing neutrons that cause further fissions; characterised by the
  multiplication factor $k$. (\chref{ch:reactions})
\item[cloud chamber] A detector in which charged particles leave visible trails of droplets in a
  supersaturated vapour. (\chref{ch:decay})
\item[Coulomb barrier ($U_B$)] The electric potential-energy hill that two positively charged nuclei must
  overcome, or tunnel through, before they can touch. (\chref{ch:reactions})
\item[critical, subcritical, supercritical] A chain reaction with $k=1$ (steady), $k<1$ (dying out) or $k>1$
  (growing). The \emph{critical mass} is the smallest mass of fissile material that can be critical.
  (\chref{ch:reactions})
\item[curve peeling] Separating a decay curve that is the sum of two exponentials by fitting the long-lived
  tail first and subtracting it. (\chref{ch:decaylaw})
\item[daughter] The nucleus produced by a decay. (\chref{ch:decay})
\item[dating (radioactive)] Finding the age of a closed sample from a ratio of parent and daughter numbers,
  $t=\Thalf\log_2(1+N_Y/N_X)$, or from its remaining activity. (\chref{ch:nucapps})
\item[de Broglie wavelength] The wavelength $h/p$ associated with any particle of momentum $p$.
  (\chref{ch:reactions})
\item[decay constant ($\lambda$)] The probability per unit time that a given nucleus decays; unit
  \si{s^{-1}}. (\chref{ch:decaylaw})
\item[decay series] A chain of successive decays from a long-lived heavy parent to a stable end; the four
  natural series are labelled $4n$, $4n+1$, $4n+2$ and $4n+3$ by their mass numbers. (\chref{ch:decay})
\item[delayed neutrons] Fission neutrons emitted seconds after the fission by decaying fragments; they make
  reactors controllable. (\chref{ch:reactions})
\item[deterministic and stochastic effects] Radiation harm that appears only above a threshold dose and grows
  in severity with dose (burns, radiation sickness), and harm whose probability, not severity, grows with
  dose (chiefly cancer). (\chref{ch:nucapps})
\item[distance of closest approach ($d$)] The nearest a charged projectile fired straight at a nucleus gets
  before the Coulomb repulsion turns it back. (\chref{ch:nucleus})
\item[effective half-life ($T_\text{eff}$)] The half-life of a substance removed by physical decay and
  biological excretion together, $1/T_\text{eff}=1/T_\text{phys}+1/T_\text{bio}$. (\chref{ch:nucapps})
\item[electron capture (EC)] A nucleus absorbs one of its own inner (K-shell) electrons, turning a proton into
  a neutron and emitting a neutrino. (\chref{ch:decay})
\item[electronvolt (\si{\electronvolt})] The energy gained by an electron crossing a potential difference of
  \SI{1}{V}, \SI{1.602e-19}{J}; nuclear energies are usually quoted in \si{\MeV}. (\chref{ch:nucleus})
\item[endoergic, exoergic] A reaction that absorbs kinetic energy ($Q<0$) or releases it ($Q>0$).
  (\chref{ch:reactions})
\item[endpoint energy] The maximum kinetic energy in a $\beta$ spectrum, equal to $Q$ for a decay to the
  ground state; the average electron energy is only about $Q/3$. (\chref{ch:decay})
\item[enrichment] Raising the fraction of uranium-235 in uranium above its natural $0.72\%$.
  (\chref{ch:reactions})
\item[equivalent dose ($H$)] Absorbed dose weighted for biological effect, $H=w_RD$; unit sievert
  (\si{\sievert}). (\chref{ch:nucapps})
\item[fission] The splitting of a heavy nucleus into two medium-mass fragments, releasing about
  \SI{200}{\MeV}. (\chref{ch:reactions})
\item[fusion] The joining of light nuclei into a heavier one, releasing energy because $\Eb/A$ increases.
  (\chref{ch:reactions})
\item[gamma ray ($\gamma$)] A high-energy photon emitted when a nucleus drops from an excited state to a lower
  one. (\chref{ch:decay})
\item[Geiger--M\"uller tube] A gas-filled detector in which each ionising particle triggers an avalanche of
  ions, giving one countable pulse. (\chref{ch:decay})
\item[half-life ($\Thalf$)] The time for half of the nuclei in a sample to decay, $\Thalf=\ln2/\lambda$; the
  \emph{partial half-life} $T_i=\ln2/\lambda_i$ of a channel is what it would be if that channel acted alone.
  (\chref{ch:decaylaw})
\item[half-value layer ($x_{1/2}$)] The absorber thickness that halves the intensity of a narrow gamma beam,
  $x_{1/2}=\ln2/\mu$. (\chref{ch:nucapps})
\item[impact parameter ($b$)] The sideways distance between a projectile's line of approach and the centre
  of the nucleus it passes. (\chref{ch:nucleus})
\item[internal conversion] An excited nucleus gives its energy directly to an inner atomic electron, which is
  ejected instead of a photon. (\chref{ch:decay})
\item[ionisation, ion pair] Knocking an electron off an atom, leaving a positive ion and a free electron (the
  pair). (\chref{ch:decay})
\item[isobar, isotone, isotope] Nuclides with the same mass number $A$; with the same neutron number $A-Z$; with
  the same atomic number $Z$. (\chref{ch:nucleus})
\item[isomer] A nuclear excited state that lives unusually long, for example technetium-99m.
  (\chref{ch:decay})
\item[K shell] The innermost electron shell of an atom. (\chref{ch:decay})
\item[Lawson criterion] The minimum product of plasma density, temperature and confinement time for a
  fusion plasma to give useful power. (\chref{ch:reactions})
\item[lepton, electron lepton number ($L_e$)] Particles that do not feel the strong force (here the electron and
  the neutrino). $L_e=+1$ for $e^-$ and $\nu_e$, $-1$ for $e^+$ and $\bar\nu_e$, $0$ for nucleons and photons;
  total $L_e$ is conserved. (\chref{ch:decay})
\item[linear attenuation coefficient ($\mu$)] The probability per unit thickness that a photon interacts in an
  absorber, so that $I=I_0e^{-\mu x}$; $1/\mu$ is the photon's mean free path. (\chref{ch:nucapps})
\item[liquid-drop model, semi-empirical mass formula] The picture of a nucleus as a drop of incompressible
  nuclear matter, and the formula for $\Eb$ built from it (volume, surface, Coulomb and asymmetry terms).
  (\chref{ch:nucleus})
\item[mass defect ($\Delta m$)] The amount by which a nucleus is lighter than its separated nucleons.
  (\chref{ch:nucleus})
\item[mass number ($A$)] The total number of nucleons in a nucleus, protons plus neutrons.
  (\chref{ch:nucleus})
\item[mean life ($\tau$)] The average lifetime of a nucleus, $\tau=1/\lambda=1.443\,\Thalf$.
  (\chref{ch:decaylaw})
\item[moderator] A material of light nuclei (water, heavy water, graphite) that slows fission neutrons by
  collisions. (\chref{ch:reactions})
\item[Monte Carlo simulation] A calculation that imitates a chance process with random numbers, for example
  giving each nucleus a probability $p=\lambda\Delta t$ of decaying in each time step. (\chref{ch:nucapps})
\item[multiplication factor ($k$)] The average number of neutrons from one fission that cause another fission.
  (\chref{ch:reactions})
\item[neutrino ($\nu_e$)] A neutral, nearly massless lepton ($L_e=+1$) that feels only the weak force; emitted
  with a positron in $\beta^+$ decay and in electron capture. (\chref{ch:decay})
\item[neutron number ($A-Z$)] The number of neutrons in a nucleus. This book writes it as $A-Z$ and keeps $N$
  for the number of undecayed nuclei. (\chref{ch:nucleus})
\item[nucleon] A proton or a neutron. (\chref{ch:nucleus})
\item[nuclide] A particular nucleus with definite $Z$ and $A$, written $\nuc{A}{Z}{X}$ or $\iso{A}{X}$.
  (\chref{ch:nucleus})
\item[pair production] The creation of an electron and a positron by a photon of more than
  $2m_ec^2=\SI{1.022}{\MeV}$ passing close to a nucleus, which takes up momentum. (\chref{ch:decay})
\item[parent] The nucleus that decays. (\chref{ch:decay})
\item[Pauli exclusion principle] No two identical particles of this kind (electrons, protons or neutrons) can
  occupy the same quantum state. (\chref{ch:nucleus})
\item[photon] A quantum of electromagnetic radiation: no charge, no mass, speed $c$. (\chref{ch:decay})
\item[positron ($e^+$)] The antiparticle of the electron: mass \SI{0.511}{\MeV}$/c^2$, charge $+e$.
  (\chref{ch:decay})
\item[positron emission tomography (PET)] Medical imaging that locates a $\beta^+$-emitting tracer from the
  pairs of back-to-back \SI{511}{keV} photons produced when its positrons annihilate. (\chref{ch:nucapps})
\item[prompt criticality] A chain reaction that is supercritical on prompt neutrons alone, $k(1-0.0065)>1$
  for uranium-235, so its power grows on the prompt time scale. (\chref{ch:reactions})
\item[proton--proton chain] The sequence of fusion reactions by which the Sun converts hydrogen into
  helium-4. (\chref{ch:reactions})
\item[$Q$-value] The rest energy converted into kinetic energy in a decay or reaction; positive when energy is
  released. (\chref{ch:decay}, \chref{ch:reactions})
\item[radiation weighting factor ($w_R$)] The dimensionless factor that converts absorbed dose to equivalent
  dose: $1$ for $\gamma$ rays, X-rays and $\beta$ particles, $2$ for protons, $20$ for $\alpha$ particles.
  (\chref{ch:nucapps})
\item[radioactivity] The spontaneous emission of particles or photons by unstable nuclei.
  (\chref{ch:decay})
\item[radiocarbon dating] Dating organic material from its remaining carbon-14 activity per gram of carbon,
  compared with the living value of about $13.8$ decays per minute per gram. (\chref{ch:nucapps})
\item[radiogenic] Made by radioactive decay, as is the lead in a zircon crystal. (\chref{ch:nucapps})
\item[saturation activity] The steady activity, equal to the production rate $\Prod$, approached when an isotope
  is made at a constant rate for many half-lives. (\chref{ch:decaylaw})
\item[secular and transient equilibrium] Steady ratios in a decay chain when the parent outlives the daughter
  greatly (equal activities, $\lambda_XN_X=\lambda_YN_Y$) or only moderately (the daughter's activity a fixed
  multiple of the parent's). (\chref{ch:decaylaw})
\item[separation energy ($S_n$, $S_p$)] The energy needed to remove one neutron or one proton from a nucleus.
  (\chref{ch:nucleus})
\item[specific activity] Activity per unit mass, in \si{\becquerel.g^{-1}}. (\chref{ch:decaylaw})
\item[strong nuclear force] The short-range attraction between nucleons that holds nuclei together.
  (\chref{ch:nucleus})
\item[threshold energy ($K_\text{th}$)] The minimum projectile kinetic energy for an endoergic reaction,
  $K_\text{th}=|Q|(1+m_a/m_X)$. (\chref{ch:reactions})
\item[tracer] A radioactive isotope followed through a body or a process by its radiation; chemically it behaves
  exactly like the stable isotope. (\chref{ch:nucapps})
\item[tunnelling (quantum)] The passage of a quantum particle through an energy barrier higher than its
  energy; it makes $\alpha$ decay and stellar fusion possible. (\chref{ch:decay}, \chref{ch:reactions})
\item[uncertainty principle] A particle confined within a distance $\Delta x$ has a momentum spread of at least
  about $\hbar/\Delta x$. (\chref{ch:decay})
\item[valley of stability] The band of stable nuclides on a chart of neutron number $A-Z$ against atomic number
  $Z$. (\chref{ch:decay})
\item[weak nuclear force] The interaction that changes neutrons into protons and vice versa; responsible for
  beta decay. (\chref{ch:decay})
\end{description}}
